\documentclass[12pt]{article}
% Préambule commun à tous les documents extraits de « La Revue CAPES »
\usepackage[T1]{fontenc}
\usepackage[utf8]{inputenc}
\usepackage{lmodern}
\usepackage{amsmath,amssymb,amsthm,amsfonts,mathrsfs}
\usepackage{setspace}
\usepackage{listings}
\usepackage{pgfplots}
\pgfplotsset{compat=1.18}
\usepackage{longtable}
\usepackage{adjustbox}
\usepackage{lettrine}
\usepackage{tkz-tab}
\usepackage{changepage}
\usepackage{pdflscape}
\usepackage{graphicx}
\usepackage{tikz}
\usepackage{xcolor}
\usepackage[a4paper,top=2cm,bottom=2.5cm,left=2.5cm,right=2cm]{geometry}
\usepackage{fancyhdr}
\usepackage{draftwatermark}
\SetWatermarkText{La RevueCapes}
\SetWatermarkLightness{0.9}
\SetWatermarkScale{2}
\usepackage{hyperref}
\hypersetup{colorlinks=true,linkcolor=blue!50!black,urlcolor=blue!50!black}

\definecolor{revue}{RGB}{20,60,120}

\pagestyle{fancy}
\fancyhf{}
\renewcommand{\headrulewidth}{0.4pt}
\setlength{\headheight}{15pt}

% \entete{Matière}{Titre}{Type de document}
\newcommand{\entete}[3]{%
  \fancyhead[L]{\small\textsc{La Revue CAPES} -- #1}%
  \fancyhead[R]{\small #2}%
  \fancyfoot[C]{\thepage}%
  \thispagestyle{plain}%
  \begin{center}
    {\color{revue}\rule{\textwidth}{1.2pt}}\\[4pt]
    {\small\textsc{Concours directs CAP-CEG / CAPES -- Burkina Faso}}\\[6pt]
    {\LARGE\bfseries\color{revue} #2}\\[6pt]
    {\large\itshape #1 -- #3}\\[2pt]
    {\color{revue}\rule{\textwidth}{1.2pt}}
  \end{center}
  \vspace{0.5cm}%
}

% Encadré pour les remarques ajoutées dans les corrigés
\newcommand{\remarque}[1]{\par\medskip\noindent\fcolorbox{revue}{revue!6}{\parbox{\dimexpr\linewidth-2\fboxsep-2\fboxrule}{\textbf{\color{revue}Remarque.} #1}}\par\medskip}


\begin{document}
\entete{Mathématiques}{Corrigé détaillé du sujet type n°3}{Correction}
\begin{spacing}{1.5}

\textbf{Exercice 1}

Soit $f :\mathbb{R}^2\longrightarrow \mathbb{R}$ la fonction ainsi définie par:
 $f(x,y)=2x^2+2y^2+2xy-x-y$
 
1. Calculons les dérivées partielles premières de $f$.
 
 $\frac{\partial f}{\partial x}(x,y)=4x+2y-1$
 
 $\frac{\partial f}{\partial y}(x,y)=4y+2x-1$
 
 
2. Déterminons l'unique point critique $A$ de $f$.

$\begin{cases}
 4x+2y-1=0(1)\\
 4y+2x-1=0(2)
 \end{cases}$
 
 $(1)-(2)\Longrightarrow x=y (3)$
 
 $(3) dans (1)\Longrightarrow x=\frac{1}{6}=y$
 
 \textbf{$A=(\frac{1}{6},\frac{1}{6})$}
 
 3. Calculons les dérivées partielles secondes de $f$.
 
 $\frac{\partial^2f}{\partial x^2}(x,y)=4$, $\frac{\partial^2 f}{\partial y^2}(x,y)=4$, $\frac{\partial^2 f}{\partial x\partial y}(x,y)=\frac{\partial^2 f}{\partial y\partial x}(x,y)=2$
 
 4. montrons que $f$ présente en $A$ un extremum en précisant sa nature.

La matrice hessienne de $f$ est donnée par :

$H=\begin{pmatrix}
4&2\\
2&4
\end{pmatrix}$

$\det H=\begin{vmatrix}
4&2\\
2&4
\end{vmatrix}=16-4=12>0$ et $\frac{\partial^2 f}{\partial x^2}(x,y)=4>0$ alors le seul point critique $A$ est un minimum local.

\medskip\textbf{Exercice 2}

1. Calculons $\Delta_1$ en donnant le résultat sous forme factorisée.

$\Delta_1= \begin{vmatrix}
 1 & 1 & 1\\
 b & c & d\\
 b^2 & c^2 & d^2
 \end{vmatrix}\begin{array}{c}
 c_2\longleftarrow c_2-c_1\\
 c_3\longleftarrow c_3-c_1
 \end{array}=\begin{vmatrix}
 1 & 0 & 0\\
 b & c-b & d-b\\
 b^2 & c^2-b^2 & d^2-b^2
 \end{vmatrix}=(c-b)(d-b)\begin{vmatrix}
 1 & 0 & 0\\
 b & 1 & 1\\
 b^2 & c+b & d+b
 \end{vmatrix}$
 
 $=(c-b)(d-b)\begin{vmatrix}
 1 & 1\\
 c+b & d+b
 \end{vmatrix}=(c-b)(d-b)(d+b-c-b)=(c-b)(d-b)(d-c)$
 
2. Montrons que $\Delta_2 =(b-a)(c-a)(d-a)\Delta_1$. 
 
 $\Delta_2=\begin{vmatrix}
 1 & 1 & 1 & 1 \\
 a & b & c & d \\
 a^2 & b^2 & c^2 & d^2 \\
 a^3 & b^3 & c^3 & d^3
\end{vmatrix} c_i\longleftarrow c_i-c_1$ avec $i=2,3,4 \quad\Delta_2=\begin{vmatrix}
 1 & 0 & 0 & 0 \\
 a & b-a & c-a & d-a \\
 a^2 & b^2-a^2 & c^2-a^2 & d^2-a^2 \\
 a^3 & b^3-a^3 & c^3-a^3 & d^3-a^3
\end{vmatrix}$

$=(b-a)(c-a)(d-a)\begin{vmatrix}
1 & 1 & 1 \\ 
b+a & c+a & d+a \\ 
b^2+ab+a^2 & c^2+ac+a^2 & d^2+ad+a^2
\end{vmatrix}\begin{array}{c}
L_2\longleftarrow L_2-aL_1\\
L_3\longleftarrow L_3-aL_2
\end{array}$

$\Delta_2=(b-a)(c-a)(d-a)\begin{vmatrix}
1 & 1 & 1 \\ 
b & c & d \\ 
b^2 & c^2 & d^2
\end{vmatrix}=\Delta_2 =(b-a)(c-a)(d-a)\Delta_1$.

3. Déduisons $\Delta_2$ sous forme factorisée.  

$\Delta_2 =(b-a)(c-a)(d-a)\Delta_1=(b-a)(c-a)(d-a)(c-b)(d-b)(d-c)$.

4. Le rang de $\Delta_2$.

Si les $a,b,c,d$ sont distinct deux à deux, alors $\Delta_2\neq 0$, donc de $rang(\Delta_2)=4$, sinon  $1\leq rang(\Delta_2)<4$:


 
 
 
\medskip\textbf{Exercice 3}

1. Déterminons le développement limité à l'ordre 3 au voisinage de 0 de $f$.
 
 On sait que : $e^u=1+\frac{u}{1!}+\frac{u^2}{2!}+\frac{u^3}{3!}+o(u^3)$ et $\sin(u)=u-\frac{u^3}{3!}+\circ(u^3)$ 
 
 Alors $f(x) = 2e^{-2x}sin(x)=2(1-2x+\frac{(-2x)^2}{2}+\frac{(-2x)^3}{6} +\circ(x^3))(x-\frac{x^3}{3!}+\circ(x^3))=2x-4x+\frac{11}{3}x^3+\circ(x^3)$
 
 2. Déduisons

$\lim_{x \to 0} \frac{2e^{-2x} \sin(x)}{x} = \lim_{x \to 0} \left( 2 - 4x + \frac{11}{3}x^2 + o(x^2) \right) = 2$

\medskip\textbf{Exercice 4}

1. Deux évènements sont dit incompatibles si les deux ne peuvent pas être réalisés en même temps, c'est-à-dire $A\cap B = \emptyset$.

2. Compatibilité de $A$ et $B$

On a:

 $\Omega(A\cap B)=\{(2,4);(4,2)\}$

 $A\cap B \neq \emptyset$ alors $A$ et $B$ ne sont pas incompatibles.
 
 3. Deux évènements $A$ et $B$ sont dit indépendants si $P(A\cap B)=P(A)\times P(B)$
 
 4. Calculons $P(A)$ et $P(B)$

$\Omega(A)=\{(1,5);(2,4);(3,3);(5,1);(4,2);(3,3)\}$

$P(A)=\frac{6}{36}=\frac{1}{6}$

$P(B)=\frac{6+6+6}{36}=\frac{18}{36}=\frac{1}{2}$

5. Calculons $P(A\cap B)$ et $P(A\bigtriangleup B)$

$P(A\cap B)=\frac{2}{36}=\frac{1}{18}$

$P(A\bigtriangleup B)=P(A)+P(B)-2P(A\cap B)$

$P(A\bigtriangleup B)=\frac{1}{6}+\frac{1}{2}-2\frac{1}{18}=\frac{10}{18}=\frac{5}{9}$

6. $P(A)\times P(B)=\frac{1}{6}\times \frac{1}{2}=\frac{1}{12}\neq P(A\cap B)$ alors $A$ et $B$ ne sont pas indépendants.

\end{spacing}
\end{document}
